\subsection{Наведение по метке}
\label{sec:python-visual-nav}

Ошибка из прошлого урока говорит, насколько метка смещена от центра кадра. \textbf{Обратная связь} превращает её в движение: измерили --- немного сдвинулись --- измерили снова. Собирают такой регулятор по шагам, проверяя каждый.

\begin{enumerate}
    \item \textbf{Измерить без движения.} Несколько кадров подряд печатаем ошибку, дрон висит. Даже у неподвижного дрона числа немного «дрожат» --- это шум измерения.
    \item \textbf{Знак поправки.} Ошибку кадра нужно связать с осью сцены. У нижней камеры симулятора при курсе 0 право кадра --- это $+X$ сцены, верх кадра --- $+Y$. Метка правее центра ($e_x > 0$) лежит по $+X$ от дрона, поэтому $v_x > 0$. Ось $y$ кадра направлена вниз: метка ниже центра ($e_y > 0$) лежит по $-Y$, поэтому у $v_y$ знак противоположный. Знак проверяют коротким ограниченным движением: ошибка после него должна уменьшиться.
    \item \textbf{Короткая команда.} Двигаются импульсом: несколько повторов \verb|set_manual_speed| по 0.1~с, затем явная нулевая скорость и новый кадр. Долгое движение по старой ошибке уводит цель из кадра.
    \item \textbf{Коэффициент и предел.} Скорость пропорциональна ошибке: $v_x = k \cdot e_x / 640$, где деление на ширину кадра делает ошибку безразмерной, а $k$ переводит её в м/с. Чем ближе метка к центру, тем медленнее дрон --- это P-регулятор. Скорость ограничивают: \verb|max(-0.3, min(0.3, v))|.
    \item \textbf{Допуск.} Внутри мёртвой зоны в несколько пикселей вокруг центра команда нулевая: иначе дрон бесконечно исправляет шум. «По центру» засчитывается, когда метка в допуске несколько кадров подряд.
    \item \textbf{Потеря метки.} Если метки нет, старая ошибка больше ничего не значит: скорость обнуляют сразу, а ожидание возвращения метки ограничивают числом кадров.
    \item \textbf{Снижение и посадка.} Центрирование и снижение --- отдельные этапы: снижаются только после подтверждённого центра, на тех же X и Y, затем снова центрируются и садятся.
\end{enumerate}

\begin{codefile}[python]{python/aruco_examples/12_01_land_on_marker.py}
\end{codefile}

\paragraph{Что происходит.} \verb|observation| возвращает ошибку метки или \verb|None|. \verb|align| --- регулятор по двум осям: при потере метки он останавливает дрон и через 10 кадров без метки сдаётся, в допуске --- держит нулевую скорость и считает кадры подряд, иначе --- задаёт ограниченную скорость, пропорциональную ошибке. Основная программа наводится на высоте 3~м, снижается до 1.5~м над измеренными X и Y, наводится ещё раз и только после этого сажает дрон. Если какой-то этап не удался, программа говорит об этом и сажает дрон на месте --- посадку «по цели» при этом успешной не объявляет.

\begin{rulebox}[Знаки и курс]
Связь «право кадра $= +X$, верх $= +Y$» верна для нижней камеры при курсе 0. После поворота корпуса оси кадра поворачиваются вместе с ним, и знаки поправок нужно пересчитать (или задавать скорость командой \verb|set_manual_speed_body_fixed| в осях дрона). У камеры настоящего дрона своё направление, поэтому знак каждой оси проверяют коротким опытом заново.
\end{rulebox}

\begin{taskbox}[Знак поправки]
Возьмите из примера \verb|observation|. Над стартом на высоте 3~м измерьте ошибку метки 12, затем одну секунду двигайтесь по оси X со скоростью 0.2~м/с (повторы по 0.1~с), явно остановитесь и измерьте снова. Знак скорости выберите так, чтобы ошибка уменьшилась, и проверьте обратный знак отдельным запуском.
\expected{при $v_x = +0.2$ значение $e_x$ уменьшается (метка 12 стоит правее центра), при $v_x = -0.2$ --- растёт; после нулевой команды дрон висит.}
\mission{Python}{91--93}
\end{taskbox}

\begin{taskbox}[Посадка по ArUco]
Запустите пример, затем поменяйте его параметры по одному: коэффициент \verb|K| 2 $\to$ 6, допуск \verb|DEAD_ZONE| 8 $\to$ 0, знак у \verb|vy|. Для каждого варианта запишите, что происходит. Верните исходные значения и проверьте ветку потери: укажите \verb|TARGET_ID = 249|, которого на карте нет.
\expected{с исходными параметрами дрон центрируется над меткой 12, снижается, садится на неё и печатает точку посадки около $(0.6, 0.3)$; с большим \texttt{K} он проскакивает центр, без допуска дольше «дрожит» около него, с неверным знаком метка уходит из кадра; с ID 249 программа печатает «Метка потеряна» и сажает дрон на месте.}
\mission{Python}{94--99}
\end{taskbox}

\subsection{Следование за меткой (только на дроне)}

Посадка выше использует ошибку в пикселях. Если метка стоит впереди, а не внизу, а нужно держать её в центре кадра и на заданном расстоянии, добавляют расстояние из \verb|tvec| (раздел~\ref{sec:python-aruco-pose}). Сначала --- следование одним поворотом, без расстояния:
\begin{enumerate}
    \item Найти центр метки на кадре (среднее арифметическое углов).
    \item Найти центр кадра: половины ширины и высоты из \texttt{frame.shape}.
    \item Если $X_{marker} > X_{center}$, метка справа $\rightarrow$ поворачиваем вправо; если меньше --- влево.
    \item Скорость поворота пропорциональна ошибке $e_x$ и ограничена, как \verb|align| из прошлого урока.
\end{enumerate}

\begin{codefile}[python]{python/aruco_examples/12_follow_aruco.py}
\end{codefile}

Этот пример --- заготовка следования для настоящего дрона. Поток обработки ищет маркер (словарь \texttt{DICT\_4X4\_50}: маркер печатают из этого словаря; симулятор такие метки не найдёт, его метки нарисованы словарём \texttt{DICT\_6X6\_250}), рисует точку в его центре и через \texttt{cv2.solvePnP} вычисляет смещение маркера относительно камеры. Метод \texttt{control\_loop} пока ничего не отправляет дрону: дрон не взлетает, а команды движения закомментированы. Дописать регулятор по пунктам списка выше --- ваша задача; устроен он так же, как \verb|align| из прошлого урока, только ошибка $e_x$ управляет скоростью поворота \verb|yaw_rate|.

\begin{taskbox}[Посадка на базу]
Только на дроне. Дрон висит на высоте 1~м. Если он видит маркер ID~1, он 1~с летит вперёд со скоростью 0.2~м/с; если ID~2 --- назад; если ID~0 --- садится. Каждое решение принимается по метке, подтверждённой тремя кадрами подряд; без меток скорость нулевая.
\expected{показ маркера 1 сдвигает дрон вперёд примерно на 20~см, маркера 2 --- назад, маркер 0 сажает его; короткое мелькание маркера в кадре дрон не двигает.}
\end{taskbox}
